\section{Zero-current switching quasi-resonant converters (ZCS-QRC)}

\subsection{Ctrl+F keywords: ZCS, QRC, quasi-resonant, half-wave, full-wave, resonant switch, hard switching, switching loss, \(\omega_0\), \(Z_0\), \(\alpha\), \(\Phi\), \(g\), linear state, resonant state}

Hard switching: \(v\) and \(i\) overlap at commutations \(\Rightarrow\) switching loss. ZCS: \(L_r\) in series with the switch forces \(i_{\mathrm{SW}}\) through zero so the device can turn off at zero current; turn-on is also soft because \(i_{Lr}\) ramps from zero.

\subsection{Resonant switch}
Half-wave: series diode, \(i_{Lr}\) cannot reverse. Full-wave: antiparallel diode, \(i_{Lr}\) can go negative.

\[
\omega_0=\frac{1}{\sqrt{L_r C_r}},\qquad f_0=\frac{\omega_0}{2\pi},\qquad Z_0=\sqrt{\frac{L_r}{C_r}}.
\]
Bulk \(L_F\), \(C_F\) are large: load port is a current sink \(I_o\) (buck) or source \(I_{\mathrm{in}}\) (boost).

\subsection{Unified gain function}
\[
g=\frac{f_s}{2\pi f_0}\left(\frac{\Phi}{2}+\alpha+\frac{1-\cos\alpha}{\Phi}\right),\qquad
F(\Phi)=\frac{\Phi}{2}+\alpha+\frac{1-\cos\alpha}{\Phi}.
\]
Full-wave: \(F(\Phi)\approx 2\pi\) for \(0<\Phi<1\), hence \(g\approx f_s/f_0\) and the gain is almost load-independent. Half-wave: \(F(\Phi)\) is not flat; \(M\) depends on \(R\).

PWM \(\leftrightarrow\) ZCS (replace \(D\) by \(g\)):
\begin{center}
\begin{tabular}{@{}llll@{}}
\toprule
 & PWM & ZCS-QRC & ZVS-QRC\\
\midrule
Buck & \(D\) & \(g\) & \(1-g\)\\
Boost & \(1/(1-D)\) & \(1/(1-g)\) & \(1/g\)\\
Buck-boost / Cuk & \(D/(1-D)\) & \(g/(1-g)\) & \((1-g)/g\)\\
\bottomrule
\end{tabular}
\end{center}

\subsection{ZCS-QR buck: four states}
Linear \([t_0,t_1]\): diode still freewheels, \(L_r\) sees \(V_{\mathrm{in}}\):
\[
i_{Lr}=\frac{V_{\mathrm{in}}}{L_r}(t-t_0),\qquad T_{d1}=\frac{L_r I_o}{V_{\mathrm{in}}},\quad i_{Lr}(t_1)=I_o.
\]
Resonant \([t_1,t_2]\):
\[
i_{Lr}=I_o+\frac{V_{\mathrm{in}}}{Z_0}\sin\omega_0(t-t_1),\qquad
v_{Cr}=V_{\mathrm{in}}\bigl(1-\cos\omega_0(t-t_1)\bigr).
\]
Ends at \(i_{Lr}=0\):
\[
T_{d2}=\frac{\alpha}{\omega_0},\qquad \alpha=\sin^{-1}\Bigl(-\frac{Z_0 I_o}{V_{\mathrm{in}}}\Bigr),
\]
\[
\pi<\alpha<\tfrac32\pi\ \text{(half-wave)},\qquad \tfrac32\pi<\alpha<2\pi\ \text{(full-wave)}.
\]
\[
v_{Cr}(t_2)=V_{\mathrm{in}}(1-\cos\alpha).
\]
Recovering \([t_2,t_3]\): \(C_r\) discharges at \(I_o\):
\[
v_{Cr}=V_{\mathrm{in}}(1-\cos\alpha)-\frac{I_o}{C_r}(t-t_2),\qquad T_{d3}=\frac{C_r V_{\mathrm{in}}(1-\cos\alpha)}{I_o}.
\]
Freewheel \([t_3,t_4]\): \(T_{d4}=T_s-T_{d1}-T_{d2}-T_{d3}\).

\subsubsection{ZCS condition (buck)}
Peak resonant current must reach zero:
\[
\frac{V_{\mathrm{in}}}{Z_0}\ge I_o \quad\Leftrightarrow\quad \Phi=\frac{Z_0 I_o}{V_{\mathrm{in}}}\le 1.
\]
Turn-off window: half-wave, after \(i_{Lr}\) hits zero and before switch voltage goes positive; full-wave, while \(i_{Lr}<0\) (antiparallel diode conducting).

\subsubsection{ZCS-QR buck conversion ratio}
\[
M=\frac{V_o}{V_{\mathrm{in}}}=\frac{f_s}{2\pi f_0}\left(\frac{\Phi}{2}+\alpha+\frac{1-\cos\alpha}{\Phi}\right)=g,\qquad
\Phi=\frac{Z_0 I_o}{V_{\mathrm{in}}}=\frac{M}{R_n},\quad R_n=\frac{R}{Z_0}.
\]
Recursive in \(\Phi\) because \(\alpha=\alpha(\Phi)\). Solve numerically unless full-wave \(g\approx f_s/f_0\) is accepted.

\subsection{ZCS-QR boost}
Linear: \(L_r\) sees \(V_o\),
\[
i_{Lr}=\frac{V_o}{L_r}(t-t_0),\qquad T_{d1}=\frac{L_r I_{\mathrm{in}}}{V_o}.
\]
Resonant:
\[
i_{Lr}=I_{\mathrm{in}}+\frac{V_o}{Z_0}\sin\omega_0(t-t_1),\qquad v_{Cr}=V_o\cos\omega_0(t-t_1),
\]
\[
\alpha=\sin^{-1}\Bigl(-\frac{Z_0 I_{\mathrm{in}}}{V_o}\Bigr),\qquad T_{d2}=\frac{\alpha}{\omega_0}.
\]
Recovering: \(C_r\) charges at \(I_{\mathrm{in}}\) until \(v_{Cr}=V_o\),
\[
T_{d3}=\frac{C_r V_o(1-\cos\alpha)}{I_{\mathrm{in}}}.
\]
ZCS:
\[
\frac{V_o}{Z_0}\ge I_{\mathrm{in}}.
\]
Gain:
\[
M=\frac{V_o}{V_{\mathrm{in}}}=\frac{1}{1-g},\qquad \Phi=\frac{Z_0 I_{\mathrm{in}}}{V_o}.
\]
\(M\) cannot exceed \(R_n\) or ZCS is lost. MOSFET source can share ground (implementation plus).

\subsection{ZCS-QR buck-boost}
Equivalent current sink is \(I_{\mathrm{in}}+I_o\). ZCS:
\[
\frac{V_{\mathrm{in}}+V_o}{Z_0}\ge I_{\mathrm{in}}+I_o.
\]
\[
M=\frac{g}{1-g},\qquad \Phi=\frac{Z_0(I_{\mathrm{in}}+I_o)}{V_{\mathrm{in}}+V_o}.
\]

\subsection{Exam recipe: full-wave ZCS-QR buck at ZCS boundary}
\textbf{Use when:} given \(V_s\), \(V_o\), \(f_s\), \(I_o\); ``operated at boundary''; find \(L_r\), \(C_r\); optional \(I_{Lr,\mathrm{rms}}\) loss.

\textbf{Steps:}
\begin{enumerate}
  \item Boundary: \(Z_0=V_s/I_o\) (\(\Phi=1\)).
  \item Full-wave approximation: \(M\approx f_s/f_0\) \(\Rightarrow\) \(f_0\approx f_s/M=f_s V_s/V_o\).
  \item \(Z_0=\sqrt{L_r/C_r}\), \(\omega_0=1/\sqrt{L_r C_r}=2\pi f_0\) \(\Rightarrow\) \(L_r=Z_0/\omega_0\), \(C_r=1/(\omega_0 Z_0)\).
  \item For ESR loss, approximate \(i_{Lr}\) (linear ramp + half/full sine) and \(P=I_{Lr,\mathrm{rms}}^2 R_{\mathrm{esr}}\).
\end{enumerate}

\textbf{Fast checks:} \(Z_0=V/I\) in ohms; \(f_0>f_s\) because buck \(M<1\).

\textbf{Common traps:} using ZVS condition \(Z I_o>V_{\mathrm{in}}\); designing half-wave \(\alpha\) while the problem says full-wave; forgetting \(f_0=\omega_0/(2\pi)\).

\subsection{Recognition rules}
ZCS waveform: switch current is a sinusoid that returns to zero; switch voltage is a pulse. Half-wave current never negative. Control variable is \(f_s\), not PWM \(D\).

\subsection{Common traps}
\(\Phi\) definition changes with topology (normalise to the resonant voltage/current pair of that converter). Half-wave is load-dependent; full-wave is not (approximately). Junction-capacitor loss still exists in ZCS (motivation for ZVS).

\subsection{One-line checks}
Buck ZCS: \(M\) rises with \(f_s/f_0\). Boost ZCS: \(M=1/(1-g)>1\). ZCS requires \(\Phi\le 1\).
